\let\negmedspace\undefined
\let\negthickspace\undefined
\documentclass[journal]{IEEEtran}
\usepackage[a5paper, margin=10mm, onecolumn]{geometry}
\usepackage{lmodern} % Ensure lmodern is loaded for pdflatex
\usepackage{tfrupee} % Include tfrupee package

\setlength{\headheight}{1cm} % Set the height of the header box
\setlength{\headsep}{0mm}     % Set the distance between the header box and the top of the text

\usepackage{gvv-book}
\usepackage{gvv}
\usepackage{cite}
\usepackage{amsmath,amssymb,amsfonts,amsthm}
\usepackage{algorithmic}
\usepackage{graphicx}
\graphicspath{{./figs/}}
\usepackage{textcomp}
\usepackage{xcolor}
\usepackage{txfonts}
\usepackage{listings}
\usepackage{enumitem}
\usepackage{mathtools}
\usepackage{gensymb}
\usepackage{comment}
\usepackage[breaklinks=true]{hyperref}
\usepackage{tkz-euclide}
\usepackage{listings}
\usepackage{gvv}
\def\inputGnumericTable{}

\usepackage[latin1]{inputenc}
\usepackage{color}
\usepackage{array}

\usepackage{longtable}
\usepackage{calc}
\usepackage{multirow}

\usepackage{hhline}
\usepackage{ifthen}

\usepackage{lscape}
\usepackage{circuitikz}

% Custom command for conjugate
\newcommand{\conj}[1]{{#1}^c}

\begin{document}
	
	\bibliographystyle{IEEEtran}
	\vspace{3cm}
	
	\title{12.367}
	\author{EE25BTECH11042 - Nipun Dasari}
	\maketitle
	
	\renewcommand{\thefigure}{\theenumi}
	\renewcommand{\thetable}{\theenumi}
	\setlength{\intextsep}{10pt}
	
	\numberwithin{equation}{enumi}
	\numberwithin{figure}{enumi}
	\renewcommand{\thetable}{\theenumi}
	
	\textbf{Question}:\\
	The eigen-values of a real symmetric matrix are always
	\begin{enumerate}[label=\alph*)]
		\item positive
		\item imaginary
		\item real
		\item complex conjugate pairs
	\end{enumerate}
	
	\solution
	
	The correct option is (c), the eigenvalues of a real symmetric matrix are always **real**.
	
	\textbf{Proof without Hermitian Conjugate}:
	
	Let $\vec{A}$ be a real symmetric matrix. This means all its entries are real numbers and $\vec{A} = \vec{A}^\top$.
	
	Let $\lambda$ be an eigenvalue of $\vec{A}$ and $\vec{v}$ be its corresponding eigenvector. Note that $\lambda$ and the components of $\vec{v}$ could be complex. The eigenvalue equation is:
	\begin{align}
		\vec{A}\vec{v} = \lambda\vec{v}
	\end{align}
	If we take the complex conjugate of this entire equation, we get:
	\begin{align}
		(\vec{A}\vec{v})^c = (\lambda\vec{v})^c \\
		\conj{\vec{A}}\conj{\vec{v}} = \conj{\lambda}\conj{\vec{v}}
	\end{align}
	Since $\vec{A}$ is a **real** matrix, its complex conjugate is itself ($\conj{\vec{A}} = \vec{A}$). So, the equation becomes:
	\begin{align}
		\vec{A}\conj{\vec{v}} = \conj{\lambda}\conj{\vec{v}} \label{eq:conj}
	\end{align}
	Now, left-multiply the original eigenvalue equation by $(\conj{\vec{v}})^\top$:
	\begin{align}
		(\conj{\vec{v}})^\top\vec{A}\vec{v} = (\conj{\vec{v}})^\top(\lambda\vec{v}) = \lambda((\conj{\vec{v}})^\top\vec{v}) \label{eq:first}
	\end{align}
	Next, take the transpose of both sides of equation (\ref{eq:conj}):
	\begin{align}
		(\vec{A}\conj{\vec{v}})^\top = (\conj{\lambda}\conj{\vec{v}})^\top \\
		(\conj{\vec{v}})^\top\vec{A}^\top = \conj{\lambda}(\conj{\vec{v}})^\top
	\end{align}
	Since $\vec{A}$ is a **symmetric** matrix, $\vec{A}^\top = \vec{A}$. The equation is:
	\begin{align}
		(\conj{\vec{v}})^\top\vec{A} = \conj{\lambda}(\conj{\vec{v}})^\top
	\end{align}
	Now, right-multiply this equation by $\vec{v}$:
	\begin{align}
		(\conj{\vec{v}})^\top\vec{A}\vec{v} = \conj{\lambda}((\conj{\vec{v}})^\top\vec{v}) \label{eq:second}
	\end{align}
	By comparing equations (\ref{eq:first}) and (\ref{eq:second}), we see that:
	\begin{align}
		\lambda((\conj{\vec{v}})^\top\vec{v}) = \conj{\lambda}((\conj{\vec{v}})^\top\vec{v}) \\
		(\lambda - \conj{\lambda})((\conj{\vec{v}})^\top\vec{v}) = 0
	\end{align}
	The term $(\conj{\vec{v}})^\top\vec{v}$ is the dot product of the vector $\conj{\vec{v}}$ with $\vec{v}$, which is the squared norm of $\vec{v}$.
	\begin{align}
		(\conj{\vec{v}})^\top\vec{v} = \sum_{i} \conj{v_i}v_i = \sum_{i} |v_i|^2 = \norm{\vec{v}}^2
	\end{align}
	Since $\vec{v}$ is an eigenvector, it cannot be the zero vector, so its norm $\norm{\vec{v}}^2$ must be a positive real number.
	
	Given that $\norm{\vec{v}}^2 \neq 0$, the only way for the equation $(\lambda - \conj{\lambda})\norm{\vec{v}}^2 = 0$ to hold is if:
	\begin{align}
		\lambda - \conj{\lambda} = 0 \implies \lambda = \conj{\lambda}
	\end{align}
	A complex number is equal to its own conjugate only if its imaginary part is zero. Therefore, $\lambda$ must be a real number.
	
\end{document}